% 01 Apr: clean up

% Principle, Not Impossible Demand [editorial title]

\subsection{Installment 8: “Folkebladet,” May 01, 1918}

A principle is commonly a foundational proposition. It means beginning, and it is applied to the first propositions or chief rules which one follows within a given domain. It is “that ground with whose validity the rest stands or falls.”

It is, in other words, the unity of a way of thought from which a multiplicity of expressions of life are derived and explained, an original unity to which thought may be led back. As such it may be what one might call, using a foreign term, an ‘ideal aim’. But in the first place, and essentially, it is a resting-point, fertilizing seed for the whole generation of thought, a point of support, an explanatory ground for the whole of the powers of development in, and a foundation for, the moral, religious, material, personal, and social life.

A Christian principle is not beyond our powers. It is a foundation upon which to stand, something that supports, that bears. Therefore Sverdrup could write—and it was among the last things he wrote (see facsimile in S. S. III.)—“A false principle is like a failing foundation.”

Dogmatics has its principles. They are treated in what is called the doctrine of principles in dogmatics, or fundamental doctrine, which constitutes almost a third of this discipline (for example, Kirms’s, v. Oettingen’s). Lønning’s little doctrine of faith, which in my school days was used at Augsburg, contained no treatment of principles, for which reason nothing further was said about it in “class.”

Ethics also has its principles. They are found in the doctrine of principles in ethics, which likewise constitutes about a third of the discipline of ethics. Since nothing in this discipline was taught at Augsburg, nothing further was said there about the principles of ethics either.

Church polity also has certain principles which are derived from a fundamental principle. The principles of church polity were not discussed to any extent in “class.” Instruction in this area came, in theory, chiefly through the church papers, and practically through the church conflict.

Thus: three groups of principles.

1. What are the dogmatic principles?
   One speaks of two great Reformation principles, which at the same time are the chief foundational propositions for all Protestant dogmatic thought: the formal principle and the material principle, first set forth by Twesten in 1829.

The formal principle is the foundational proposition concerning Scripture as the sole decisive norm of Christianity. It is the epistemological point of departure, the source to which we must go. It is the principle of Scripture.

The material principle is the foundational proposition of justification by faith alone. It is the religious principle for a dogmatic evaluation of Christianity. It may also be designated as the Kingdom of God, a revealed act of divine grace, realized in Christ’s righteousness.

Dogmatics has to do with knowledge. On the one hand, it has a direction and an aim. On the other hand, it has a regulative principle, something that guides. The Danish dogmatician T. R. Vang says that “the direction and goal of Paul’s understanding” are “God’s thoughts concerning human beings, the content of his will, his plan of salvation and its fulfillment.” And then he asks: “But what now is the principle that guides him in this his energetic thinking?” “He has expressed it in the words that return again and again: in Christ.”

2. If one asks: What is the principle of ethics? then the answers differ. In ancient ethics it was pride. In Abelard’s ethics (the Middle Ages) it was conscience. Reinhard (ca. 1800) held that the principle of ethics was love (against Kant’s principle, pure morality). Goethe said: “I bow before him (Christ) as the divine revelation of the highest principle of morality;” but in Goethe’s eyes Christianity was surely only a revelation of natural religion. He also wanted to worship the sun, “for it is likewise a revelation of the highest!” De Wette made faith the principle of ethics, Harleß regeneration, Darwin evolution.

According to J. T. Beck, the principle of ethics is “Christ living and working in us through the Word.” Gisle Johnson: “The Christian faith, as it shows itself working through love.” Luthardt: “Communion with God in Christ, which in faith becomes a personal reality, and in life shows itself working in love.” N. Seeberg: “Man’s living communion with God in order to serve God with a view to the aim of the eternal ideal.” Grützmacher: “The source and principle of theological ethics is the Christian character.” Oftedal, Gunnersen, and Sverdrup: “Christ’s liberating power in a distinctly personal Christian life.”

3. If we were to seek a highest principle for polity, a fundamental thought from which, for example, the principles of Guiding Principles are derived, it would have to be the universal priesthood, which Kahnis regarded as the third principle of the Reformation. But whereas Kahnis called it the Church principle and applied it to the invisible Church, S. and O. applied it to the visible body of Christ, the congregation, and called it the congregational principle.

Guiding Principles are foundational propositions for church polity derived from one of the chief principles of the Reformation, the universal priesthood.

It is to be regretted that S. and O. did not have time to give expression to their principles in a rounded, concluding work. With such a work they could have made—for the time being—a contribution in practical theology which they could scarcely have made by any other work. In the area of polity they represented, for more than a generation—and with all diligence—a direction which must be described as new within Lutheran theology. Here their strength was greater than it was in dogmatics, exegesis, or history. Sverdrup was a gifted teacher in dogmatics, and he was well at home in the Old Testament. But did he represent any new direction in these disciplines? If so, wherein did it consist? Oftedal had wide-ranging interests, but did he represent any new direction in New Testament exegesis or in church history? In these areas, compared with the field of polity, they were essentially eclectics.

I do not say that this was less important than that which they contributed to polity. But in the history of theology, in the area of the concept of the Church, their exposition of church polity will stand as their distinct contribution. Even the little that has been included in “Veiledning” can, in its significance, be placed on a par with Høfling’s well-known work, Fundamental Propositions for Evangelical Lutheran Church Polity (year 1850, 225 pages, 3 editions in two years, also translated into Norwegian). Høfling’s book was described as ‘epoch-making’ and ‘of enduring significance.’ It is not too much to say that the contribution of S. and O. (even in Veiledning) just as fully deserves the same descriptions. In several respects the contribution of S. and O. signifies a great advance.

Høfling’s book ought to be able to contribute a little toward setting some “Doubting Thomas” among us on the right track, since it both says and proves that church polity goes deeper than merely external forms and various trifles, about which those who know least of the discipline never grow tired of talking. In the preface Høfling says:

\begin{quote}
It is not single forms of Evangelical-Lutheran church polity, but properly its principles and foundational propositions, concerning which in our day, alongside great confusion, there prevails a great difference of opinion, and concerning which in various places a controversy has arisen that threatens to become serious and could even give rise to a lamentable schism, if it should continue to be conducted in the same manner as hitherto and if it should not succeed in being settled by way of conviction. A manifestly Catholicizing hierarchical basic conception seeks to assert itself within the sphere of our Church, even if not as exclusively legitimate, then at any rate first of all as equally legitimate, in that it supposes itself to display a quite particular fidelity to the Church and its confession. The less now the author can regard the important meaning of our confession, or the true consequences of our Church principle with respect to this point, as doubtful, the more must he, in the interest of that Church and that confession to which he belongs with all his heart, feel himself called to restatement and rebuttal. And his highest wish is now that he may be so fortunate as to set forth the distinctive principles and foundational propositions of Evangelical-Lutheran church polity in a manner convincing to all. When he believes that he gives here one thing or another that is new, this applies not so much to practical results as to understanding of principle, development, systematic presentation, and justification…
\end{quote}

So much for Høfling.

Høfling represents the low-church standpoint in the ecclesiastical controversy in Germany around 1850. The leaders of the high-church party who attacked him were men such as Stahl, Löhe, Lechler, Huschke, Kliefoth, Dieckhoff, Münchmeyer—well-known names in church history. These men laid great weight on the Church’s legal determinations, regarded the ministry as a divine institution, and held semi-Catholic notions concerning the administration of the sacraments, ordination, and so forth (Lønning became their faithful disciple).

Høfling’s writing became a liberating word. The tendency he combated suffered a great defeat. There are now exceedingly few in the Lutheran lands of Europe who represent this tendency. Its successors came to America. Here that tendency has become a dominant force.
